\section{From GraphRAG to Epistract}
\label{sec:comparison}

Epistract grew directly from my earlier GraphRAG prototype~\cite{edge2024graphrag}, which demonstrated that LLM-augmented knowledge graphs could accelerate biomedical literature review. That prototype, however, relied on embedding similarity to define edges---a design that, while useful for clustering, left scientists unable to distinguish a side-effect relationship from a mechanistic one. Table~\ref{tab:evolution} summarises how each architectural decision evolved between the two systems, and Figure~\ref{fig:evolution} illustrates the progression visually.

\begin{table}[t]
\centering
\caption{Architectural evolution from GraphRAG to Epistract.}
\label{tab:evolution}
\begin{tabular}{@{}lll@{}}
\toprule
\textbf{Dimension} & \textbf{GraphRAG (Prequel)} & \textbf{Epistract (Current)} \\
\midrule
Relationship basis   & Embedding similarity        & LLM comprehension \\
Schema               & Generic NER                 & 17 entity types, 30 relations, 40+ ontologies \\
Relation typing      & Untyped edges               & Typed, directional (\textsc{inhibits}, \textsc{causes}, etc.) \\
Evidence             & None                        & Source passage + confidence score \\
Validation           & None                        & RDKit, Biopython, NCT/CAS patterns \\
Development tool     & Cursor + Gemini 2.5 Pro     & Claude Code \\
Delivery             & Scripts, manual setup        & One-command plugin install \\
Scale tested         & 6 articles                  & 77 articles across 5 domains \\
Cross-domain         & Single domain               & 5 domains, zero retraining \\
Community labels     & Numbered                    & Semantic heuristic labeling \\
Reproducibility      & Code on GitHub              & Plugin + committed test artifacts \\
\bottomrule
\end{tabular}
\end{table}

\paragraph{Typed relations recover scientific semantics.}
In GraphRAG, an edge between \emph{nivolumab} and \emph{colitis} is an undirected similarity link---the scientist cannot determine whether it represents a side effect, a treatment target, or a textual co-mention. In Epistract, the same relationship is captured as \texttt{nivolumab \textrightarrow{} CAUSES \textrightarrow{} colitis} with a confidence score of 0.85 and an accompanying source passage. A scientist can therefore query ``what does nivolumab cause?'' and receive a direct, verifiable answer rather than a ranked list of nearest neighbours in an embedding space.

\paragraph{Evidence traceability supports regulatory-grade review.}
Every relation in an Epistract graph carries both a confidence score and the exact text passage from which it was extracted. This design enables scientists to verify any edge back to the source literature---a requirement for pharmaceutical decision-making that similarity-based graphs cannot satisfy. When a drug-safety reviewer asks why two entities are connected, the system returns not a cosine distance but a quotation and a provenance chain.

\paragraph{Molecular validation closes the loop between text and chemistry.}
GraphRAG treats SMILES strings and amino acid sequences as opaque text tokens, indistinguishable from any other string. Epistract validates them with RDKit and Biopython, computing structural properties (InChIKey, molecular weight, Lipinski parameters) and sequence properties (GC content, isoelectric point). A SMILES string in an Epistract graph is therefore a verified molecular structure, not an LLM approximation. This validation step catches hallucinated or malformed chemical identifiers before they propagate into downstream analyses.

\paragraph{The epistemological distinction.}
Taken together, these changes reveal a fundamental difference in what each system captures. GraphRAG asks \emph{what entities appear near each other in the text?} Epistract asks \emph{what does the text say about how these entities relate?} The former captures statistical co-occurrence; the latter captures scientific meaning. For drug discovery, where the difference between ``inhibits'' and ``activates'' determines whether a compound is a therapeutic candidate or a safety risk, this distinction is not academic---it is operational. Epistract does not dismiss the GraphRAG approach; it builds on its proof of concept and replaces proximity with comprehension.
